% ============================================
% Assistant Professor Cover Letter – Texas A&M University (Macroeconomics)
% ============================================

\documentclass[11pt,letterpaper]{letter}

\usepackage{geometry}
\usepackage{microtype}
\usepackage[hidelinks]{hyperref}

\geometry{left=25mm,right=25mm,top=25mm,bottom=25mm}

\signature{Decory Edwards}
\address{
Department of Economics \\
Johns Hopkins University \\
Baltimore, MD 21218 \\
\texttt{dedwar65@jh.edu}
}

\begin{document}

\begin{letter}{Search Committee \\
Department of Economics \\
Texas A\&M University \\
College Station, TX \\ USA}

\opening{Dear Members of the Search Committee,}

I am writing to apply for the tenure-track Assistant Professor position in Macroeconomics in the Department of Economics at Texas A\&M University. I am a Ph.D.\ candidate in Economics at Johns Hopkins University, advised by Professors Christopher D.~Carroll, Nicholas W.~Papageorge, and Stelios Fourakis, and I expect to complete my degree in May~2026. My research lies at the intersection of macroeconomics, heterogeneous-agent models, and wealth inequality, with an emphasis on how heterogeneity in returns and income risk shapes aggregate dynamics and policy transmission.

My job market paper estimates a distribution of heterogeneous returns to wealth using micro data from the Survey of Consumer Finances and embeds these estimates into a structural heterogeneous-agent macroeconomic model. I show that persistent return heterogeneity plays a central role in generating observed wealth concentration and has important implications for macroeconomic policy, particularly for the distributional effects of monetary and fiscal interventions. A related project uses longitudinal data from the Health and Retirement Study to study how trust and beliefs interact with portfolio choices and realized returns, providing a behavioral channel linking micro-level heterogeneity to macroeconomic outcomes. Together, these projects form a coherent research agenda focused on distributional macroeconomics, grounded in micro data and quantitative modeling.

Texas A\&M’s Economics Department is an excellent environment for this work. The department’s strengths in macroeconomics, applied microeconomics, and econometrics—along with its emphasis on policy-relevant research—closely match my research interests and training. I am particularly excited by the opportunity to contribute to a large, research-active department with strong graduate and undergraduate programs and to engage with colleagues whose work spans theory, empirical analysis, and public policy. Texas A\&M’s role as a flagship public research university also aligns well with my interest in research that speaks to broad economic questions with real-world relevance.

\textbf{Teaching.} I have experience teaching and supporting courses in \textit{Principles of Macroeconomics}, \textit{Intermediate Macroeconomic Theory}, and upper-level electives in macroeconomic policy. My teaching philosophy emphasizes clarity, economic intuition, and the use of data to connect theory to real-world outcomes. At Texas A\&M, I would be well prepared to teach undergraduate and graduate macroeconomics, including principles, intermediate macro, and field courses in macroeconomic policy or distributional macroeconomics. I am also eager to contribute to graduate training and to mentor students in empirical and quantitative research.

Beyond academic fit, Texas A\&M is an especially attractive institution for me given its strong academic culture, commitment to excellence in research and teaching, and location within a region that I know well and value highly. I would be enthusiastic about building a long-term academic career at Texas A\&M and contributing to the department’s mission in research, teaching, and service.

Thank you very much for your time and consideration. I would welcome the opportunity to discuss my application further.

\closing{Sincerely,}

\end{letter}
\end{document}
